\documentclass[11pt,a4paper]{article}

% ============================================================
% Language, typography, and layout
% ============================================================
\usepackage[utf8]{inputenc}
\usepackage[T1]{fontenc}
\usepackage[spanish,es-tabla]{babel}
\usepackage{lmodern}
\usepackage{microtype}
\usepackage[a4paper,margin=2.35cm]{geometry}

% ============================================================
% Mathematics
% ============================================================
\usepackage{amsmath}
\usepackage{amssymb}
\usepackage{bm}

% ============================================================
% Figures, tables, and references
% ============================================================
\usepackage{graphicx}
\usepackage{booktabs}
\usepackage{float}
\usepackage{subcaption}
\usepackage{hyperref}
\usepackage{xcolor}
\usepackage{siunitx}

\hypersetup{
  colorlinks=true,
  linkcolor=blue!55!black,
  citecolor=blue!55!black,
  urlcolor=blue!55!black,
  pdftitle={Estructura de bandas fotonicas de metamateriales dispersivos},
  pdfauthor={Hermitian photonic bands reproduction},
  pdfsubject={Hermitian eigenvalue formulation for dispersive photonic crystals},
  pdfkeywords={photonic bands, dispersive metamaterials, Hermitian eigenvalue problem, Drude model}
}

% ============================================================
% Macros
% ============================================================
\newcommand{\ii}{\mathrm{i}}
\newcommand{\dd}{\mathrm{d}}
\newcommand{\epsinf}{\varepsilon_\infty}
\newcommand{\omegap}{\omega_p}
\newcommand{\omegao}{\omega_0}
\newcommand{\vect}[1]{\bm{#1}}
\newcommand{\Hhat}{\widehat{H}}

% ============================================================
% Title
% ============================================================
\title{
  \textbf{Estructura de bandas fotónicas de metamateriales dispersivos}\\[0.35em]
  \large Reproducción numérica de la formulación Hermitiana de Raman y Fan
}

\author{
  Reproducción computacional basada en \textit{Phys. Rev. Lett.} 104, 087401 (2010)
}

\date{\today}

\begin{document}

\maketitle

\begin{abstract}
Se presenta una reproducción numérica de la formulación Hermitiana para el
cálculo de estructuras de bandas fotónicas en metamateriales dispersivos
propuesta por Raman y Fan. El sistema considerado es una red periódica
bidimensional de barras cuadradas plasmónicas en aire, descritas por un modelo
de Drude sin pérdidas para el cálculo de bandas y por un término disipativo
débil para el análisis modal de pérdidas. La introducción de campos auxiliares
permite transformar el problema dispersivo en un problema lineal de autovalores.
En el caso sin pérdidas, dicho problema es Hermitiano y conduce a frecuencias
reales y modos ortogonales con una norma energética. En el caso disipativo, se
comparan las pérdidas obtenidas mediante diagonalización no-Hermitiana con la
predicción perturbativa de primer orden. Los resultados reproducen las
características principales reportadas para las bandas TE: aparición de bandas
planas asociadas a modos plasmónicos superficiales, perfiles de campo
localizados y buena concordancia entre pérdida exacta y perturbativa para
$\gamma = 0.01\omegap$.
\end{abstract}

\tableofcontents
\newpage

% ============================================================
\section{Introducción}
% ============================================================

Los cristales fotónicos y metamateriales periódicos permiten controlar la
propagación de ondas electromagnéticas mediante la ingeniería espacial de sus
propiedades ópticas. En materiales dieléctricos no dispersivos, el cálculo de
bandas fotónicas se formula como un problema Hermitiano de autovalores. Esta
estructura matemática garantiza frecuencias reales, ortogonalidad modal y una
base adecuada para desarrollar análisis espectrales.

En metamateriales plasmónicos y medios dispersivos, la permitividad depende de
la frecuencia. Por ello, la ecuación estándar para el campo electromagnético ya
no define directamente un problema lineal de autovalores. La formulación de
Raman y Fan resuelve esta dificultad introduciendo campos auxiliares asociados
a los grados de libertad materiales. Con estos campos, la dispersión se absorbe
en un sistema aumentado con coeficientes independientes de la frecuencia.

El objetivo de este trabajo es reproducir el caso bidimensional de barras
cuadradas plasmónicas en aire, en la polarización TE, y verificar dos resultados
principales:

\begin{enumerate}
  \item la estructura de bandas sin pérdidas obtenida mediante un problema
        Hermitiano;
  \item la comparación entre pérdidas modales exactas y pérdidas estimadas por
        teoría de perturbaciones.
\end{enumerate}

% ============================================================
\section{Modelo dispersivo}
% ============================================================

El material dispersivo se modela mediante una respuesta de Lorentz:

\begin{equation}
  \varepsilon(\omega)
  =
  \epsinf
  \left(
    1 +
    \frac{\omegap^2}
    {
      \omegao^2 - \omega^2 + \ii\gamma\omega
    }
  \right),
  \label{eq:lorentz}
\end{equation}

donde $\epsinf$ es la permitividad de alta frecuencia, $\omegap$ es la
frecuencia de plasma, $\omegao$ es la frecuencia natural del oscilador y
$\gamma$ representa amortiguamiento.

El caso metálico considerado corresponde al límite Drude:

\begin{equation}
  \omegao = 0.
\end{equation}

Para la reproducción numérica se emplean unidades normalizadas tales que

\begin{equation}
  a = 1,
  \qquad
  c = 1,
  \qquad
  \mu_0 = 1,
  \qquad
  \omegap = 1,
\end{equation}

donde $a$ es la constante de red.

En el metal se toma

\begin{equation}
  \epsinf = 1,
  \qquad
  \omegap = 1,
  \qquad
  \omegao = 0.
\end{equation}

En aire se toma

\begin{equation}
  \epsinf = 1,
  \qquad
  \omegap = 0,
\end{equation}

de modo que no existe un grado dispersivo activo en la región no metálica.

% ============================================================
\section{Campos auxiliares}
% ============================================================

La respuesta material se describe mediante la polarización $\vect{P}$, cuya
dinámica satisface

\begin{equation}
  \frac{\partial^2 \vect{P}}{\partial t^2}
  +
  \gamma
  \frac{\partial \vect{P}}{\partial t}
  +
  \omegao^2 \vect{P}
  =
  \omegap^2 \epsinf \vect{E}.
  \label{eq:polarization}
\end{equation}

Se introduce la velocidad de polarización

\begin{equation}
  \vect{V}
  =
  \frac{\partial \vect{P}}{\partial t}.
  \label{eq:v_def}
\end{equation}

Entonces las ecuaciones acopladas de Maxwell y del oscilador material pueden
escribirse como

\begin{align}
  \frac{\partial \vect{H}}{\partial t}
  &=
  -\frac{1}{\mu_0}\nabla\times\vect{E},
  \label{eq:h_time}
  \\
  \frac{\partial \vect{E}}{\partial t}
  &=
  \frac{1}{\epsinf}
  \left(
    \nabla\times\vect{H} - \vect{V}
  \right),
  \label{eq:e_time}
  \\
  \frac{\partial \vect{P}}{\partial t}
  &=
  \vect{V},
  \label{eq:p_time}
  \\
  \frac{\partial \vect{V}}{\partial t}
  &=
  \omegap^2\epsinf\vect{E}
  -
  \omegao^2\vect{P}
  -
  \gamma\vect{V}.
  \label{eq:v_time}
\end{align}

En el límite Drude, $\omegao=0$, el bloque asociado a $\vect{P}$ puede
eliminarse del sistema reducido.

% ============================================================
\section{Formulación Hermitiana}
% ============================================================

En régimen armónico, el sistema aumentado puede escribirse como

\begin{equation}
  \omega A x = Bx,
  \label{eq:generalized_evp}
\end{equation}

donde $x$ contiene los campos electromagnéticos y auxiliares. En el caso sin
pérdidas, $\gamma = 0$, la matriz $B$ es Hermitiana y $A$ es Hermitiana definida
positiva.

Se define el cambio de variable

\begin{equation}
  y = A^{1/2}x.
  \label{eq:y}
\end{equation}

Entonces el problema se transforma en

\begin{equation}
  \omega y = \Hhat y,
  \qquad
  \Hhat = A^{-1/2}BA^{-1/2}.
  \label{eq:hhat_problem}
\end{equation}

Como $\Hhat$ es Hermitiana, sus autovalores son reales y sus autovectores
pueden escogerse ortonormales. En variables físicas, la ortogonalidad modal se
expresa como

\begin{equation}
  x_m^\dagger A x_n = \delta_{mn}.
  \label{eq:orthogonality}
\end{equation}

Esta norma coincide con la energía total del modo. Para el caso Drude, la
densidad de energía sin pérdidas puede escribirse como

\begin{equation}
  W_0
  =
  \frac{1}{2}\epsinf |\vect{E}|^2
  +
  \frac{1}{2}\mu_0 |\vect{H}|^2
  +
  \frac{1}{2\epsinf\omegap^2}
  |\vect{V}|^2.
  \label{eq:energy}
\end{equation}

% ============================================================
\section{Polarización TE en dos dimensiones}
% ============================================================

El sistema es bidimensional y uniforme en la dirección $z$. Para la polarización
TE, el campo magnético tiene componente $H_z$, mientras que el campo eléctrico
y el campo auxiliar están en el plano:

\begin{equation}
  x_{\mathrm{TE}}
  =
  \left(
    H_z,\,
    E_x,\,
    E_y,\,
    V_x,\,
    V_y
  \right)^T.
  \label{eq:te_state}
\end{equation}

Esta es la polarización empleada para la reproducción de la estructura de
bandas de barras cuadradas plasmónicas.

% ============================================================
\section{Geometría y parámetros numéricos}
% ============================================================

La celda unitaria contiene una barra cuadrada metálica centrada en aire. El
lado de la barra es

\begin{equation}
  s = 0.25a.
  \label{eq:side}
\end{equation}

La ruta de Bloch considerada es $\Gamma \rightarrow X$:

\begin{equation}
  k_x \in \left[0,\frac{\pi}{a}\right],
  \qquad
  k_y = 0.
\end{equation}

Las frecuencias se reportan en la forma normalizada

\begin{equation}
  \frac{\omega a}{2\pi c}.
  \label{eq:normalized_frequency}
\end{equation}

Los parámetros principales se resumen en la Tabla~\ref{tab:parameters}.

\begin{table}[H]
  \centering
  \caption{Parámetros usados en la reproducción numérica.}
  \label{tab:parameters}
  \begin{tabular}{lll}
    \toprule
    Magnitud & Valor & Descripción \\
    \midrule
    $a$ & $1$ & constante de red \\
    $c$ & $1$ & velocidad de la luz normalizada \\
    $\mu_0$ & $1$ & permeabilidad normalizada \\
    $s$ & $0.25a$ & lado de la barra metálica \\
    $\epsinf$ & $1$ & metal y aire \\
    $\omegap$ & $1$ & metal \\
    $\omegap$ & $0$ & aire \\
    $\omegao$ & $0$ & modelo Drude \\
    $\gamma$ & $0$ & estructura de bandas sin pérdidas \\
    $\gamma$ & $0.01\omegap$ & cálculo con pérdidas \\
    Resolución & $20\times20$ & discretización espacial principal \\
    Polarización & TE & campos $(H_z,E_x,E_y,V_x,V_y)$ \\
    \bottomrule
  \end{tabular}
\end{table}

% ============================================================
\section{Discretización}
% ============================================================

La celda unitaria se discretiza con una malla de Yee. Esta discretización
coloca los campos eléctricos y magnéticos en posiciones escalonadas, lo que
preserva de forma natural la estructura rotacional de las ecuaciones de
Maxwell.

La periodicidad se incorpora mediante condiciones de Bloch. Para un campo
$f(\vect{r})$,

\begin{equation}
  f(x+a,y) = e^{\ii k_x a}f(x,y),
  \qquad
  f(x,y+a) = e^{\ii k_y a}f(x,y).
  \label{eq:bloch}
\end{equation}

La propiedad discreta central es la relación adjunta entre los operadores de
rotacional eléctrico y magnético:

\begin{equation}
  C_H = C_E^\dagger.
  \label{eq:curl_adjoint}
\end{equation}

Esta relación garantiza que el operador aumentado sea Hermitiano en ausencia
de pérdidas.

% ============================================================
\section{Pérdidas modales}
% ============================================================

Cuando $\gamma > 0$, el problema deja de ser Hermitiano y las frecuencias
adquieren parte imaginaria:

\begin{equation}
  \omega = \omega' + \ii\omega''.
\end{equation}

Con la convención temporal adoptada, el decaimiento modal corresponde a
$\omega''<0$. Por ello, la pérdida se representa como $-\operatorname{Im}
(\omega)$.

La predicción perturbativa de primer orden para el caso Drude es

\begin{equation}
  \omega_1
  =
  -\frac{\ii\gamma}{2}
  \frac{
    \displaystyle
    \int_{\mathrm{metal}}
    \frac{|\vect{V}_0|^2}{\epsinf\omegap^2}
    \,\dd \vect{r}
  }{
    \displaystyle
    \int W_0\,\dd \vect{r}
  }.
  \label{eq:perturbation}
\end{equation}

Esta expresión muestra que la pérdida modal es proporcional a la fracción de
energía almacenada en el grado mecánico dispersivo.

% ============================================================
\section{Resultados}
% ============================================================

\subsection{Estructura de bandas TE}

La Figura~\ref{fig:bands} muestra la estructura de bandas TE para barras
cuadradas plasmónicas en aire. Se observan bandas relativamente planas en
rangos de frecuencia asociados a modos plasmónicos superficiales. Esta es una
de las firmas principales del sistema dispersivo: la respuesta del metal
introduce resonancias locales que generan bandas con baja dispersión.

\begin{figure}[H]
  \centering
  \includegraphics[width=0.92\textwidth]{../../results/figures/fig1_square_rods_TE_bands.png}
  \caption{Estructura de bandas TE para barras cuadradas plasmónicas en aire.
  La ruta de Bloch es $\Gamma\rightarrow X$ y la frecuencia se reporta como
  $\omega a/2\pi c$.}
  \label{fig:bands}
\end{figure}

\subsection{Perfiles de campo}

Las Figuras~\ref{fig:field_mode0} y~\ref{fig:field_mode4} muestran perfiles
espaciales representativos de los modos en el punto $\Gamma$. Se presentan
componentes del campo eléctrico y del campo auxiliar de velocidad de
polarización. La comparación entre $E_x$ y $V_x$ permite identificar qué modos
concentran mayor energía en el material dispersivo.

\begin{figure}[H]
  \centering
  \includegraphics[width=0.86\textwidth]{../../results/field_profiles/fig1_TE_mode0_Gamma.png}
  \caption{Perfil espacial del modo TE 0 en $\Gamma$, normalizado por energía
  total.}
  \label{fig:field_mode0}
\end{figure}

\begin{figure}[H]
  \centering
  \includegraphics[width=0.86\textwidth]{../../results/field_profiles/fig1_TE_mode4_Gamma.png}
  \caption{Perfil espacial de un modo TE adicional en $\Gamma$. La intensidad
  del campo auxiliar permite estimar cualitativamente la pérdida modal.}
  \label{fig:field_mode4}
\end{figure}

\subsection{Pérdida exacta y perturbativa}

La Figura~\ref{fig:loss} compara la pérdida modal obtenida por diagonalización
no-Hermitiana con la predicción de primer orden. Para
$\gamma=0.01\omegap$, ambos métodos muestran buena concordancia. La mayor
pérdida aparece en regiones espectrales donde los modos presentan mayor
carácter plasmónico y, por tanto, mayor participación del campo auxiliar
$\vect{V}$.

\begin{figure}[H]
  \centering
  \includegraphics[width=0.86\textwidth]{../../results/figures/fig2_loss_comparison.png}
  \caption{Comparación entre pérdida modal exacta y pérdida calculada mediante
  teoría de perturbaciones de primer orden para $\gamma=0.01\omegap$.}
  \label{fig:loss}
\end{figure}

\subsection{Validez perturbativa}

La Figura~\ref{fig:perturbation_gamma} muestra el error relativo entre el
método exacto y el método perturbativo para diferentes valores de $\gamma$.
Como se espera para una expansión de primer orden, la aproximación mejora al
disminuir la intensidad de la pérdida.

\begin{figure}[H]
  \centering
  \includegraphics[width=0.86\textwidth]{../../results/figures/perturbation_vs_gamma.png}
  \caption{Error relativo entre la pérdida exacta y la pérdida perturbativa
  para diferentes valores de $\gamma$.}
  \label{fig:perturbation_gamma}
\end{figure}

\subsection{Convergencia espacial}

La Figura~\ref{fig:convergence} muestra la convergencia espacial de una
frecuencia representativa al refinar la malla. La reducción del error relativo
con la resolución confirma que las cantidades calculadas convergen hacia un
valor estable al aumentar el número de puntos por celda.

\begin{figure}[H]
  \centering
  \includegraphics[width=0.78\textwidth]{../../results/figures/convergence_test.png}
  \caption{Convergencia espacial de la frecuencia fundamental respecto a la
  resolución de la malla.}
  \label{fig:convergence}
\end{figure}

% ============================================================
\section{Discusión}
% ============================================================

Los resultados reproducen los elementos físicos centrales del caso estudiado
por Raman y Fan. La formulación Hermitiana permite obtener una estructura de
bandas real en el caso sin pérdidas, aun cuando el material sea dispersivo. La
presencia de bandas planas indica modos con fuerte localización asociada a la
respuesta plasmónica del metal.

Los perfiles de campo muestran que la intensidad del campo auxiliar no es
uniforme entre modos. Esta observación es fundamental para entender la pérdida:
dos modos con frecuencias similares pueden tener pérdidas distintas si almacenan
fracciones diferentes de su energía en el grado material dispersivo.

La comparación entre el cálculo exacto no-Hermitiano y la teoría de
perturbaciones confirma que, para pérdidas débiles, la base modal lossless
proporciona una descripción cuantitativa de la disipación. En particular, la
corrección perturbativa predice correctamente que la pérdida escala con la
energía cinética del oscilador material.

% ============================================================
\section{Conclusiones}
% ============================================================

Se reprodujo numéricamente la formulación Hermitiana para bandas fotónicas en
un metamaterial dispersivo de tipo Drude. La introducción del campo auxiliar
$\vect{V}$ permite linealizar la dependencia en frecuencia y obtener un
problema Hermitiano en ausencia de pérdidas.

Para barras cuadradas plasmónicas en aire se obtuvo una estructura de bandas TE
con bandas planas asociadas a modos plasmónicos superficiales. Al introducir
una pérdida débil $\gamma=0.01\omegap$, la parte imaginaria de las frecuencias
calculada mediante diagonalización no-Hermitiana concuerda con la predicción
perturbativa de primer orden.

Estos resultados confirman que la formulación Hermitiana no solo permite
calcular bandas de medios dispersivos, sino que también proporciona una base
ortonormal físicamente significativa para analizar pérdidas modales.

% ============================================================
\section*{Referencia principal}
% ============================================================

El desarrollo teórico reproducido en este documento se basa en el trabajo de
Raman y Fan~\cite{Raman2010}.

\bibliographystyle{plain}
\bibliography{bibliography}

\end{document}
